\documentclass{amsart}
% For dealing with references we use the comment environment
\usepackage{verbatim}
\newenvironment{reference}{\comment}{\endcomment}
%\newenvironment{reference}{}{}
\newenvironment{slogan}{\comment}{\endcomment}
\newenvironment{history}{\comment}{\endcomment}

% For commutative diagrams we use Xy-pic
\usepackage[all]{xy}

% We use 2cell for 2-commutative diagrams.
\xyoption{2cell}
\UseAllTwocells

% We use multicol for the list of chapters between chapters
\usepackage{multicol}

% This is generally recommended for better output
\usepackage{lmodern}
\usepackage[T1]{fontenc}

% For cross-file-references

% Package for hypertext links:
\usepackage{hyperref}

% For any local file, say "hello.tex" you want to link to please

% Theorem environments.
%
\theoremstyle{plain}
\newtheorem{theorem}[subsection]{Theorem}
\newtheorem{proposition}[subsection]{Proposition}
\newtheorem{lemma}[subsection]{Lemma}

\theoremstyle{definition}
\newtheorem{definition}[subsection]{Definition}
\newtheorem{example}[subsection]{Example}
\newtheorem{exercise}[subsection]{Exercise}
\newtheorem{situation}[subsection]{Situation}

\theoremstyle{remark}
\newtheorem{remark}[subsection]{Remark}
\newtheorem{remarks}[subsection]{Remarks}

\numberwithin{equation}{subsection}

% Macros
%
\def\lim{\mathop{\mathrm{lim}}\nolimits}
\def\colim{\mathop{\mathrm{colim}}\nolimits}
\def\Spec{\mathop{\mathrm{Spec}}}
\def\Hom{\mathop{\mathrm{Hom}}\nolimits}
\def\Ext{\mathop{\mathrm{Ext}}\nolimits}
\def\SheafHom{\mathop{\mathcal{H}\!\mathit{om}}\nolimits}
\def\SheafExt{\mathop{\mathcal{E}\!\mathit{xt}}\nolimits}
\def\Sch{\mathit{Sch}}
\def\Mor{\mathop{\mathrm{Mor}}\nolimits}
\def\Ob{\mathop{\mathrm{Ob}}\nolimits}
\def\Sh{\mathop{\mathit{Sh}}\nolimits}
\def\NL{\mathop{N\!L}\nolimits}
\def\CH{\mathop{\mathrm{CH}}\nolimits}
\def\proetale{{pro\text{-}\acute{e}tale}}
\def\etale{{\acute{e}tale}}
\def\QCoh{\mathit{QCoh}}
\def\Ker{\mathop{\mathrm{Ker}}}
\def\Im{\mathop{\mathrm{Im}}}
\def\Coker{\mathop{\mathrm{Coker}}}
\def\Coim{\mathop{\mathrm{Coim}}}

% Boxtimes
%
\DeclareMathSymbol{\boxtimes}{\mathbin}{AMSa}{"02}

%
% Macros for moduli stacks/spaces
%
\def\QCohstack{\mathcal{QC}\!\mathit{oh}}
\def\Cohstack{\mathcal{C}\!\mathit{oh}}
\def\Spacesstack{\mathcal{S}\!\mathit{paces}}
\def\Quotfunctor{\mathrm{Quot}}
\def\Hilbfunctor{\mathrm{Hilb}}
\def\Curvesstack{\mathcal{C}\!\mathit{urves}}
\def\Polarizedstack{\mathcal{P}\!\mathit{olarized}}
\def\Complexesstack{\mathcal{C}\!\mathit{omplexes}}
% \Pic is the operator that assigns to X its picard group, usage \Pic(X)
% \Picardstack_{X/B} denotes the Picard stack of X over B
% \Picardfunctor_{X/B} denotes the Picard functor of X over B
\def\Pic{\mathop{\mathrm{Pic}}\nolimits}
\def\Picardstack{\mathcal{P}\!\mathit{ic}}
\def\Picardfunctor{\mathrm{Pic}}
\def\Deformationcategory{\mathcal{D}\!\mathit{ef}}

\setlength{\emergencystretch}{3em}
\begin{document}
\title[Stacks Project, Tag 055G]{471 --- Geometric connectedness of the generic fibre spreads for finite-type morphisms over irreducible bases}
\maketitle
\noindent Copyright (C) 2005--2025 Johan de Jong. Stacks Project authors. Source: \href{https://stacks.math.columbia.edu/tag/055G}{Tag 055G}.

\noindent Editorial difficulty note (not author text). Difficulty H2: The proof first passes to an étale base change making generic components geometrically irreducible. It then spreads their closures and their pairwise incidence graph simultaneously, using flatness of the nonempty intersections to preserve the graph and connectedness over every geometric fibre..

\noindent Editorial provenance note (not author text). History: excerpt from revision \texttt{a0444\allowbreak 6e57e\allowbreak c1fbc\allowbreak 252a8\allowbreak 71afc\allowbreak ec775\allowbreak 2fb28\allowbreak 07b14}, more-morphisms.tex, lines 7943--8033. Reference presentation was adapted; original-author context and core support blocks were retained or added. The mathematical wording is unchanged. Licensed under GNU FDL 1.2 or later; no invariant sections or cover texts. See ../licenses/COPYING. Context blocks reproduce author definitions and supporting results; they are not additional counted problems.

\begin{lemma}
\label{lemma-geometrically-connected-generic-fibre}
Let $f : X \to Y$ be a morphism of schemes.
Assume
\begin{enumerate}
\item $Y$ is irreducible with generic point $\eta$,
\item $X_\eta$ is geometrically connected, and
\item $f$ is of finite type.
\end{enumerate}
Then there exists a nonempty open subscheme $V \subset Y$
such that $X_V \to V$ has geometrically connected fibres.
\end{lemma}

\begin{proof}
Choose a diagram
$$
\xymatrix{
X' \ar[d]_{f'} \ar[r]_{g'} & X_V \ar[r] \ar[d] & X \ar[d]^f \\
Y' \ar[r]^g & V \ar[r] & Y
}
$$
as in
Lemma \href{https://stacks.math.columbia.edu/tag/0551}{lemma make components generic fibre geometrically irreducible (Tag 0551)}.
Note that the generic fibre of $f'$ is geometrically connected
(for example by
Lemma \href{https://stacks.math.columbia.edu/tag/055F}{lemma base change fibres nr geometrically connected components (Tag 055F)}).
Suppose that the lemma holds for the morphism $f'$. This means that
there exists a nonempty open $W \subset Y'$ such that every fibre of
$X' \to Y'$ over $W$ is geometrically connected.
Then, as $g$ is an open morphism by
Morphisms, Lemma \href{https://stacks.math.columbia.edu/tag/03WT}{lemma etale open (Tag 03WT)}
all the fibres of $f$ at points of the nonempty open $V = g(W)$ are
geometrically connected, see
Lemma \href{https://stacks.math.columbia.edu/tag/055F}{lemma base change fibres nr geometrically connected components (Tag 055F)}.
In this way we see that we may assume that the irreducible
components of the generic fibre $X_\eta$ are geometrically irreducible.

\medskip\noindent
Let $Y'$ be the reduction of $Y$, and set $X' = Y' \times_Y X$.
Then it suffices to prove the lemma for the morphism $X' \to Y'$
(for example by
Lemma \href{https://stacks.math.columbia.edu/tag/055F}{lemma base change fibres nr geometrically connected components (Tag 055F)}
once again). Since the generic fibre of $X' \to Y'$ is the same as the
generic fibre of $X \to Y$ we see that we may assume that $Y$ is
irreducible and reduced (i.e., integral, see
Properties, Lemma \href{https://stacks.math.columbia.edu/tag/01ON}{lemma characterize integral (Tag 01ON)})
and that the irreducible
components of the generic fibre $X_\eta$ are geometrically irreducible.

\medskip\noindent
At this point suppose that
$X_\eta = X_{1, \eta} \bigcup \ldots \bigcup X_{n, \eta}$
is the decomposition of the generic fibre into
(geometrically) irreducible components.
Let $X_i$ be the closure of $X_{i, \eta}$ in $X$.
After shrinking $Y$ we may assume that
$X = \bigcup X_i$, see
Lemma \href{https://stacks.math.columbia.edu/tag/054Y}{lemma cover generic fibre neighbourhood (Tag 054Y)}.
Let $Z_{i, j} = X_i \cap X_j$.
Let
$$
\{1, \ldots, n\} \times \{1, \ldots, n\} = I \amalg J
$$
where $(i, j) \in I$ if $Z_{i, j, \eta} = \emptyset$ and
$(i, j) \in J$ if $Z_{i, j, \eta} \not = \emptyset$.
After shrinking $Y$ we may assume that $Z_{i, j} = \emptyset$
for all $(i, j) \in I$, see
Lemma \href{https://stacks.math.columbia.edu/tag/054W}{lemma empty generic fibre (Tag 054W)}.
After shrinking $Y$ we obtain that $X_{i, y}$
is geometrically irreducible for each $i$ and all $y \in Y$, see
Lemma \href{https://stacks.math.columbia.edu/tag/0559}{lemma geom irreducible generic fibre (Tag 0559)}.
After shrinking $Y$ some more we achieve the situation where
each $Z_{i, j} \to Y$ is flat and of finite presentation for
all $(i, j) \in J$, see
Morphisms, Proposition \href{https://stacks.math.columbia.edu/tag/052A}{proposition generic flatness (Tag 052A)}.
This means that $f(Z_{i, j}) \subset Y$ is open, see
Morphisms, Lemma \href{https://stacks.math.columbia.edu/tag/01UA}{lemma fppf open (Tag 01UA)}.
We claim that
$$
V  = \bigcap\nolimits_{(i, j) \in J} f(Z_{i, j})
$$
works, i.e., that $X_y$ is geometrically connected for each
$y \in V$. Namely, the fact that $X_\eta$ is connected implies that
the equivalence relation generated by the pairs in $J$ has only
one equivalence class. Now if $y \in V$ and $K \supset \kappa(y)$
is a separably closed extension, then the irreducible components
of $(X_y)_K$ are the fibres $(X_{i, y})_K$. Moreover, we see by
construction and $y \in V$ that $(X_{i, y})_K$ meets $(X_{j, y})_K$
if and only if $(i, j) \in J$. Hence the remark on equivalence classes
shows that $(X_y)_K$ is connected and we win.
\end{proof}
\end{document}
